% =====================================================================
\section{Материалы для педагога}
% =====================================================================
Этот раздел помогает встроить пособие в программу объединения. Все цифры примерные: распределение
часов зависит от возраста и подготовки группы, числа рабочих мест и наличия деталей.

\subsection*{Кому адресован курс}

Курс рассчитан на обучающихся, которые впервые знакомятся с микроконтроллерами (ориентировочно 7--11
классы). Знание языка C++ не требуется: нужные конструкции вводятся по мере необходимости и собраны
в приложении «Шпаргалка по C++». Желательны уверенная работа с компьютером и знание физики на уровне
понятий «напряжение», «ток», «сопротивление».

\subsection*{Планируемые результаты}

По окончании курса обучающийся:
\begin{itemize}
  \item соблюдает правила безопасной работы с электронными компонентами и источниками питания;
  \item читает простые принципиальные схемы, собирает их на макетной плате и проверяет мультиметром;
  \item пишет программы для ESP32-S3: работает с цифровыми и аналоговыми выводами, ШИМ, таймерами
        и прерываниями, строит программу без блокирующих задержек;
  \item подключает датчики и исполнительные устройства по интерфейсам UART, I\textsuperscript{2}C, SPI, 1-Wire;
  \item использует Wi-Fi и Bluetooth LE для обмена данными и управления устройством;
  \item находит причину неисправности по сообщениям платы, измерениям и документации;
  \item самостоятельно проектирует, собирает, документирует и защищает собственное устройство.
\end{itemize}

\subsection*{Примерное тематическое планирование (72 часа)}

Вариант рассчитан на год занятий по два академических часа в неделю.

\begin{center}\small
\begin{tabularx}{\textwidth}{@{}X l c@{}}
\toprule
Раздел & Главы пособия & Часы \\
\midrule
Введение: безопасность, макетная плата, сигналы, выводы и программа & вводные главы & 2 \\
Основы: выводы, Serial, кнопки, ШИМ, АЦП, I\textsuperscript{2}C & уроки 0--6 & 14 \\
Система: FreeRTOS, Wi-Fi, сон, SPI, память & уроки 7--11 & 10 \\
Связь и надёжность: ESP-NOW, MQTT, BLE, HTTPS, OTA, диагностика & уроки 12--14 & 6 \\
Датчики и исполнительные устройства & уроки 15--17, 23--26 & 14 \\
Сеть, USB и камера & уроки 18--20 & 6 \\
Приёмы программирования: таймеры и автоматы, биты, порядок в коде & уроки 21, 22, 28 & 6 \\
Автономное питание & урок 27 & 2 \\
Итоговый проект: разработка и защита & урок 41 & 12 \\
\midrule
\textbf{Всего} & & \textbf{72} \\
\bottomrule
\end{tabularx}
\end{center}

\textbf{Сокращённый вариант (36 часов):} вводные главы (2 ч), уроки 0--9 (20 ч), уроки 15 и 16 (4 ч),
мини-проект по упрощённой схеме урока 41 (10 ч).

Урок 8 расширен блоком HTML, CSS и JavaScript. В короткой программе достаточно запустить готовую
панель; подробный разбор и задания можно перенести в проектные часы или отдельный модуль.

Уроки не обязательно проходить подряд. После урока 6 порядок можно менять: в строке «Опирается на»
в начале каждого урока перечислено, что нужно знать заранее. Уроки 21, 22 и 28 полезно пройти раньше,
если группа быстро осваивает материал.

\subsection*{Как построить занятие}

\begin{center}\small
\begin{tabularx}{\textwidth}{@{}l c X@{}}
\toprule
Этап & Минуты & Содержание \\
\midrule
Вход в тему & 10 & вопросы «Проверь себя» прошлого урока; раздел «Сначала» нового урока \\
Объяснение & 20 & теория с рисунками пособия; разбор готового примера по строкам \\
Сборка и запуск примера & 25 & каждый собирает схему и добивается работы примера; раздел «Если не работает» \\
Самостоятельный опыт & 25 & изменение примера под новую задачу; помощь педагога по запросу \\
Итог & 10 & что получилось, типичные ошибки, задание на дом по уровню (\ding{72}--\ding{72}\ding{72}\ding{72}\ding{72}) \\
\bottomrule
\end{tabularx}
\end{center}

\subsection*{Формы контроля}

\begin{itemize}
  \item \textbf{Текущий контроль} --- вопросы «Проверь себя» (ответы в приложении) и работающий
        самостоятельный опыт на каждом занятии.
  \item \textbf{Промежуточный контроль} --- шесть заданий четырёх уровней сложности в конце уроков
        (два \ding{72}, два \ding{72}\ding{72}, одно \ding{72}\ding{72}\ding{72}, одно \ding{72}\ding{72}\ding{72}\ding{72});
        достаточно одного выполненного задания на урок, уровень обучающийся выбирает сам, но
        рекомендуйте идти по порядку и не пропускать лёгкие.
  \item \textbf{Итоговая аттестация} --- защита проекта по критериям урока 41.
\end{itemize}

\subsection*{Организационные советы}

\begin{itemize}
  \item Первое занятие начинайте с главы «Техника безопасности и рабочее место»; правила стоит повторять
        перед уроками 16, 25 и 27, где появляются моторы и аккумуляторы.
  \item Одно рабочее место --- один комплект из приложения «Компоненты для курса». Работа в парах
        допустима, но каждый должен сам собрать схему хотя бы через раз.
  \item Имя и пароль учебной сети Wi-Fi выдавайте на занятии; не используйте школьную административную
        сеть. Для уроков 8, 12, 13 и 18 удобна отдельная точка доступа \SI{2.4}{GHz}.
  \item Перед уроками 15, 17 и 18 заранее установите библиотеки на всех компьютерах: скачивание во время
        занятия отнимает много времени.
  \item Если деталей не хватает, часть опытов уроков 0--6 и 21--23 можно выполнять в симуляторе Wokwi.
  \item Урок 20 (камера) требует другой платы; его можно провести как демонстрацию.
\end{itemize}

\subsection*{Расширение по ESP32-S3: ещё 40 часов}
Базовое планирование на 72 часа выше сохраняется как отдельный маршрут: уроки 0--28 и проект 41.
Новые уроки 29--38 образуют дополнительный модуль на 40 часов: по четыре часа на тему.
Маршрут без веб-модуля 39--40 --- 112 часов. Он требует отдельных рабочих мест для проводных шин,
аудиомодулей, экранов и установленной ESP-IDF для уроков 35--36.
На 36-часовом маршруте новые темы можно использовать как демонстрации, без обещания полного освоения.

Изучайте модуль 29--38 до веб-модуля 39--40 и проекта 41; для JTAG и ULP заранее проверьте ESP-IDF и доступ к нативному USB.

\subsection*{Веб-панели и настройка: ещё 8 часов}
Урок 39 (ESPUI) --- четыре часа: готовые элементы, callbacks, очередь, два браузера и телеметрия.
Урок 40 (WiFiManager) --- четыре часа: режимы настройки и работы, окно доступа и восстановление.
Полный маршрут теперь 120 часов: базовые 72, расширение ESP32-S3 40, веб-модуль 8.
При недостатке времени минимальную ESPUI-панель можно показать после урока 8; очередь требует урока 7.
Зависимости заранее установите на все рабочие места. Для проверки нужны два браузера и отдельная
учебная сеть 2,4 ГГц; аппаратные приводы в первом опыте замените LED.
